This section covers the X-FAB XH018 implementation of Variant~A.
The technology choice and the area budget are given in the product overview
(Table~\ref{tab:asTechAreaA}).

\subsection{SRAM Macro Selection -- Variant A}
\label{subsec:sramSelectionA}

XH018 offers several SPRAM compiler families (high speed, low power, very low
power, low voltage), each with and without a \emph{bit} write-enable mask
(the ``B'' variants).
In the compiled sizes needed here, the high-speed and low-power macros have
practically the same area. The low-power family is chosen because
Variant~A runs at a low core clock
(Section~\ref{subsec:clockReqNoPipeline}), so the faster macro brings no benefit.

\begin{table}[H]
\caption{Variant A -- XH018 SPRAM Instances}
\label{tab:asSramInstancesA}
\centering
\begin{tabularx}{\textwidth}{|l |l |l |l |X|}
  \hline
  Memory & Organisation & Size & Compiler & Write enable \\
  \hline
  \hline
  I-Mem & 8192 $\times$ 32\,bit & 32\,KiB & SPRAMLP  & one \texttt{WEn} for the whole word \\
  \hline
  D-Mem & 1024 $\times$ 64\,bit & 8\,KiB  & SPRAMBLP & \texttt{WEn[63:0]}, one per bit \\
  \hline
\end{tabularx}
\end{table}

\paragraph{Why a bit-write macro for the D-Mem.}
The D-Mem must support byte, halfword, word and doubleword stores
(\texttt{SB}/\texttt{SH}/\texttt{SW}/\texttt{SD}). \texttt{asDMemCore} already
generates an 8-bit byte enable \texttt{byteEn\_i[7:0]}.
A macro with only one write enable would need either eight 8-bit-wide macros
or a read-modify-write cycle for every partial store.
The bit-write macro avoids both: each byte enable drives the eight mask bits
of its byte.

\paragraph{Why a plain macro for the I-Mem.}
The I-Mem is written only by the JTAG loader, always as complete 32-bit
words, and is read-only while the CPU runs. No write mask is needed.

\subsection{SRAM Interface -- Variant A}
\label{subsec:sramIfA}

Both macros are synchronous single-port RAMs with one clock edge per access.
All control inputs are \textbf{active low}.

\begin{table}[H]
\caption{Variant A -- SPRAM Port Mapping}
\label{tab:asSramPortsA}
\centering
\begin{tabularx}{\textwidth}{|l |l |X|}
  \hline
  SPRAM port & Direction & Connected to \\
  \hline
  \hline
  \texttt{CLK}  & in  & memory clock (rising edge) \\
  \hline
  \texttt{CEn}  & in  & low = access in this cycle; high = idle \\
  \hline
  \texttt{WEn}  & in  & low = write. I-Mem: 1\,bit. D-Mem: 64\,bit,
                        \texttt{WEn[8i+7:8i] = \{8\{\textasciitilde byteEn[i]\}\}} \\
  \hline
  \texttt{A}    & in  & word address (I-Mem 13\,bit, D-Mem 10\,bit) \\
  \hline
  \texttt{D}    & in  & write data \\
  \hline
  \texttt{Q}    & out & read data, valid after the clock edge \\
  \hline
  \texttt{RDY}  & out & test output, not used functionally \\
  \hline
\end{tabularx}
\end{table}

Unlike the high-speed family, the low-power macros have no output enable
(\texttt{OEn}); \texttt{Q} is always driven.

Without \texttt{`define XFAB\_SPRAM} the wrappers contain a behavioural
model of our own, so every simulator works without X-FAB files. The X-FAB
Verilog models are covered by the NDA and are not part of the repository.
The I-Mem macro is clocked through a multiplexer: \texttt{tck\_i} while
JTAG loads the program, \texttt{clk\_div\_s} otherwise
(Table~\ref{tab:asImemModes}).

\subsection{Pad Ring -- Variant A}
\label{subsec:padRingA}

The chip top \texttt{cc\_chip\_top} (\texttt{src/cc\_chip\_top.sv}) contains
the pad ring and the core \texttt{as\_top\_mem} (instance \texttt{core}).
Its ports are the bond pads and keep the former \texttt{as\_top\_mem} port
names, so the SDC port names hold for both levels; it is the Genus top.
The pads are instantiated through three wrappers:

\begin{itemize}
  \item \texttt{cc\_pad\_in} -- input, CMOS or Schmitt trigger
    (parameter \texttt{schmitt}),
  \item \texttt{cc\_pad\_out} -- output, tri-state cell with output enable
    (active high; tied active except for TDO),
  \item \texttt{cc\_pad\_bidir} -- bidirectional (GPIO), output enable
    active high.
\end{itemize}

As with the SRAM wrappers, \texttt{`define XFAB\_PADS} selects the
IO\_CELLS\_3V cells; the cell names come from the local path file
(\texttt{\$XFAB\_PATHS}, NDA). Without the define the wrappers are
behavioural, and \texttt{cc\_chip\_top} also serves as FPGA top with the old
pin names. The NAND-tree test chain of the pad cells (\texttt{PI}/\texttt{PO})
is not used.
The tri-state drivers moved out of the core into the pads: the GPIO driver
from \texttt{as\_gpio} into the bidirectional pads (\texttt{as\_top\_mem}
ports \texttt{gpio\_o}, \texttt{gpio\_oe\_o}, \texttt{gpio\_i},
Table~\ref{tab:asGpioPer04a}), the TDO driver from \texttt{jtag} into the
TDO pad (new port \texttt{tdo\_oe\_o}, active only in Shift-IR/Shift-DR, so
that TDO is high impedance otherwise, as required by IEEE~1149.1). The core therefore contains no internal tri-state
net. The testbenches resolve the GPIO tri-state themselves.

\begin{table}[H]
\caption{Variant A -- Signal Pads}
\label{tab:asPadsA}
\centering
\begin{tabularx}{\textwidth}{|l |l |l |X|}
  \hline
  Pin & Dir. & Pad & Remark \\
  \hline
  \hline
  \texttt{clk\_i}      & in  & CMOS input    & 100\,MHz \\
  \hline
  \texttt{rst\_i}      & in  & Schmitt input & asynchronous reset \\
  \hline
  \texttt{tck\_i}      & in  & Schmitt input & JTAG clock, $\le$\,25\,MHz \\
  \hline
  \texttt{trst\_i}     & in  & Schmitt input & TAP reset \\
  \hline
  \texttt{tms\_i}, \texttt{tdi\_i} & in & CMOS input & JTAG \\
  \hline
  \texttt{uart0\_rx\_i} & in & CMOS input    & UART \\
  \hline
  \texttt{tdo\_o}      & out & 4\,mA tri-state output & enable \texttt{tdo\_oe\_o} \\
  \hline
  \texttt{cs\_o}, \texttt{uart0\_tx\_o} & out & 4\,mA output & \\
  \hline
  \texttt{gpio\_io[7:0]} & bidi & 4\,mA bidirectional & 8 pads \\
  \hline
  \multicolumn{3}{|l|}{Signal pads} & 18 \\
  \hline
\end{tabularx}
\end{table}

\paragraph{Supply pads.}
Four supply roles with two pads each: core supply and core ground
(1.8\,V), IO supply and IO ground (3.3\,V) -- 8 pads. Together with the
18 signal pads this gives 26 pads plus 4 corner cells. Role, count and
placement are a first estimate from the library data and must be checked
against the IO\_CELLS\_3V databook (TECH-06). The supply pads are
instantiated in \texttt{cc\_chip\_top} only with \texttt{`define XFAB\_PADS};
corner and filler cells are added in place and route.

\paragraph{Pad-limited or core-limited?}
The core needs about 3.5--4.5\,mm\textsuperscript{2}
(Table~\ref{tab:asTechAreaA}), i.e.\ a core side of about 2\,mm. With the
pad ring around it the die side is about 2.5--2.7\,mm. A ring of this size
has room for more than 100 pads; 26 are needed.
A die that is just large enough for 26 pads would have a side below 1\,mm.
\textbf{Variant~A is core-limited}: the SRAM and the logic set the die size,
the pad ring is mostly filled with filler cells. Additional supply pads
cost no area.

\subsection{Open Items -- Technology}
\label{subsec:techOpenA}

\begin{table}[H]
\caption{Technology Open Items (Variant A)}
\label{tab:asTechOpenA}
\centering
\begin{tabularx}{\textwidth}{|l |X|}
  \hline
  ID & Item \\
  \hline
  \hline
  TECH-01 & \emph{Closed 2026-11-20:} \texttt{as\_imem} and
            \texttt{as\_dmem\_core} instantiate the macros with
            \texttt{`define XFAB\_SPRAM}, otherwise the behavioural models.
            Both were compared with the X-FAB Verilog models (random
            accesses, identical read data, no timing violations). \\
  \hline
  TECH-05 & \emph{Closed 2026-11-21:} two constraint modes in
            \texttt{scripts/xh018\_mmmc.tcl}: \texttt{func}
            (\texttt{as\_top\_mem.sdc} + \texttt{as\_top\_mem\_func.sdc},
            \texttt{im\_mode\_o}\,=\,0) and \texttt{jtag\_load}
            (\texttt{as\_top\_mem.sdc} + \texttt{as\_top\_mem\_jtag.sdc},
            \texttt{im\_mode\_o}\,=\,1, macro clocked by \texttt{tck\_i}).
            Setup views \texttt{av\_slow}, \texttt{av\_typ},
            \texttt{av\_jtag\_slow}; hold views \texttt{av\_fast},
            \texttt{av\_jtag\_fast}. \\
  \hline
  TECH-02 & \emph{Closed 2026-11-20:} the Genus scripts
            (\texttt{scripts/00\_rv64\_elab.tcl}, \texttt{01b\_rv64\_syn.tcl})
            target XH018 with \texttt{D\_CELLS\_HD} and \texttt{IO\_CELLS\_3V}.
            Corners slow/typ/fast in \texttt{scripts/xh018\_mmmc.tcl},
            constraints in \texttt{scripts/as\_top\_mem*.sdc}. The X-FAB file
            locations are read from a local file outside the repository
            (\texttt{\$XFAB\_PATHS}, NDA). Not yet run on the Cadence
            workstation. \\
  \hline
  TECH-03 & Decide the I-Mem size (32\,KiB or 16\,KiB) from the size of the
            target software. \\
  \hline
  TECH-04 & \emph{Closed 2026-11-21:} pad ring in \texttt{cc\_chip\_top}
            (Section~\ref{subsec:padRingA}): 18 signal pads, 8 supply pads,
            4 corners; core-limited. \\
  \hline
  TECH-06 & Check the supply pads (roles, count, placement, ESD and
            simultaneous switching) against the IO\_CELLS\_3V databook. \\

  \hline
\end{tabularx}
\end{table}
